% ============================================================
%  front/preface.tex
% ============================================================

\chapter*{Preface}
\addcontentsline{toc}{chapter}{Preface}

Every observation is incomplete.

This statement is so familiar that it rarely attracts attention.
Maps leave things out.
Photographs flatten depth.
Memories omit details.
Scientific instruments reveal some features while hiding others.
Even our most sophisticated mathematical models preserve certain
structures while discarding the rest.

Yet this commonplace fact raises a surprisingly deep question.

\medskip
\begin{center}
\itshape
If every observation is incomplete, how do we know anything at all?
\end{center}
\medskip

This book is an attempt to answer that question.

The traditional approach to science begins with objects and asks how
they behave.
The approach taken here begins with observation and asks what can be
reconstructed from it.
The shift is subtle but important.
Rather than treating observation as transparent access to reality,
we treat it as a structured projection from a richer domain into a
more limited one.

\bigskip

The central claim of this book can be expressed in ordinary language:

\principle{
  Observation is not access.\\[4pt]
  Observation is admissible projection.
}

The same claim can be expressed mathematically:

\principle{
  \(
  \Monoturn
  \xrightarrow{\;\Fdecomp\;}
  \text{Admissible Local Data}
  \xrightarrow{\;\pi\;}
  \text{Observable State}
  \),
  \quad
  \(
  \kernel(\pi\circ\Fdecomp)\neq 0
  \)
}

The meaning of this diagram is simple.
Reality, whatever its ultimate nature, is not presented directly to
an observer.
Instead, it is filtered through a process that determines which
information is accessible and which is not.
The resulting observable state contains genuine information about the
world, but it does not contain all of it.

The lost information is not necessarily destroyed.
It may simply lie beyond the reach of a particular observer,
instrument, scale, or horizon.

The question that follows is the question that motivates every
chapter of this book:

\principle{
  Given a projection with nontrivial kernel,\\[4pt]
  what global structure can be reconstructed from what remains?
}

The answer will take us through maps, photographs, memory, machine
learning, category theory, sheaf cohomology, agency, quantum
measurement, and cosmology.
These subjects may appear unrelated.
The argument of this book is that they are all instances of the same
underlying problem: each studies the relationship between what
exists, what is observed, and what can be reconstructed.

\section*{A note on mathematical status}

This book is careful to distinguish four levels of claim:

\begin{itemize}[leftmargin=2em]
  \item \statusproven{} --- statements that follow from established
    mathematics or physics by standard arguments.
  \item \textbf{[Established framework]} --- mathematical structures
    that are well-defined and self-consistent, even where specific
    theorems remain incomplete.
  \item \statusconjectured{} --- claims that are strongly motivated
    and precisely stated but not yet proved.
  \item \statusopen{} --- problems that remain genuinely open,
    stated as research targets.
\end{itemize}

The reader can therefore track exactly where the framework stands.
One of the recurring pathologies in theoretical physics is the
tendency to present conjectural mathematics as established
mathematics.
This book resists that tendency by marking its own uncertainty
explicitly.

\section*{How to read this book}

The book is organized so that the cosmological interpretation
appears only after the mathematical machinery that supports it has
been fully developed.

\textbf{Part~I} introduces the central ideas through familiar
examples that require no prior mathematics beyond basic algebra.

\textbf{Parts~II--III} develop the variational, dynamical, and
statistical-mechanical foundations.

\textbf{Part~IV} provides the functional-analytic well-posedness
theory that determines which mathematical structures are legitimate.

\textbf{Part~V} introduces the admissibility sheaf, reconstruction
theory, and the category-theoretic framework.

\textbf{Part~VI} states the major open problems explicitly, so that
the reader enters the cosmological chapters knowing exactly what has
and has not been established.

\textbf{Part~VII} applies the framework to cosmological structure
formation, redshift, and comparison with standard cosmology.

A reader primarily interested in the mathematics may read Parts~I--VI
as a self-contained treatment of admissibility, reconstruction, and
sheaf cohomology, with cosmology appearing as a worked example.
A reader primarily interested in cosmology may read Part~I and then
proceed directly to Part~VII, returning to the intervening parts for
technical support as needed.

\bigskip
\begin{flushright}
\itshape
Flyxion\\
Canada
\end{flushright}
